\begin{table}[tbp]
\centering
\caption{数据与质量总览}
\label{tab:t01_data_overview}
\footnotesize
\setlength{\tabcolsep}{3.5pt}
\begin{tabular}{lll}
\toprule
层面 & 指标 & 数值 \\
\midrule
样本 & 啤酒样本 / 风格数 & 250 / 22 \\
样本 & 各风格样本量（最小 / 中位 / 最大） & 3 / 10 / 31 \\
特征 & 化学特征列 & 231（实测化合物 226 + 派生汇总 5） \\
特征 & 感官维度 & 50 \\
缺失 & 化学缺失单元格 & 73 格（0.13\%） \\
缺失 & 感官缺失单元格 & 0 格（0.00\%） \\
检出限 & 下限哨兵值单元格 & 7514 格（13.01\%）；涉及 133 个特征 \\
变换 & 数值变换列数 & log10 198 列；原始尺度 33 列 \\
风格结构 & Kruskal–Wallis FDR < 0.05 & 化学 204/231；感官 49/50 \\
质量标记 & 样本 / 特征 / 不可用 & 7 / 85 / 0 \\
数据门 & D0–D6 通过数 & 7/7 \\
\bottomrule
\end{tabular}
\end{table}
% 注（写入图注/表注，勿直接编译）：“缺失”按原始 1.csv/4.csv 直接计数；下限哨兵值比例按 250 × 231 化学单元格计。风格检验为 Kruskal–Wallis 后 BH-FDR；质量标记的定义见 results/qc/report.json。
